% UCL Research Paper Declaration Form.
%
% Question wording is transcribed verbatim from the UCL Doctoral School form
% (ucl-research-paper-declaration-form.pdf, section lettering a-o as it appears
% in submitted theses). Nothing here is paraphrased: a form whose questions have
% been reworded is no longer evidence that the official form was completed.
%
% Section 1 is skipped on the form's own instruction ("if not yet published,
% please skip to section 2") because the manuscript is under review, not
% published. It is stated explicitly rather than left blank, so that an empty
% field cannot be read as an omission.
%
% One form covers both submissions. They are the same manuscript, under the same
% title, sent to two non-archival venues; splitting them across two forms would
% imply two independent pieces of work.
%
% The page-style group follows the pattern used by Abstract/abstract.tex, so the
% running head survives onto the form's second page.

\clearpage
{
\pagestyle{rpdf}
\chapter*{UCL Research Paper Declaration Form}
\thispagestyle{rpdf}
\addcontentsline{toc}{chapter}{UCL Research Paper Declaration Form}

\noindent\textit{Referencing the candidate's own published work(s).}

\vspace{0.4cm}

\noindent\textbf{1. For a research manuscript that has already been published
(if not yet published, please skip to section 2)}

\vspace{0.2cm}

\noindent Not applicable. The manuscript described below has been accepted at a
non-archival workshop and is under review at a second. Neither outcome
constitutes publication: a non-archival venue issues no version of record and
takes no copyright, so the work remains unpublished and section~1 is skipped as
the form directs.

\vspace{0.5cm}

\noindent\textbf{2. For a research manuscript prepared for publication but that
has not yet been published (if already published, please skip to section 3)}

\begin{description}[leftmargin=1.4em, style=nextline, itemsep=0.35em]
    \item[(k) What is the current title of the manuscript?]
    Routing Is Least Learnable Where It Is Most Valuable: Bounds on
    Representation Routing for Web Agents

    \item[(l) Has the manuscript been uploaded to a preprint server? (e.g.\
    medRxiv; if `Yes', please give a link or doi)]
    No.

    \item[(m) Where is the work intended to be published? (e.g.\ journal names)]
    The manuscript has been submitted, in the same form, to two
    \emph{non-archival} workshops. Non-archival submission does not constitute
    publication and does not transfer copyright, and neither venue issues a
    version of record.
    \begin{itemize}[leftmargin=1.2em, itemsep=0.15em, topsep=0.25em]
        \item REALM Workshop at EMNLP~2026 (OpenReview submission~192),
              submitted 6~August~2026 and \textbf{accepted 8~September~2026},
              with the camera-ready version due 14~September~2026:
              \url{https://openreview.net/forum?id=EAplLx6gCD}
        \item Workshop on Grounded and Faithful Vision-Language Models for
              Real-World Deployment (VLM4RWD) at NeurIPS~2026, submitted
              21~August~2026 and under review.
    \end{itemize}

    \item[(n) List the manuscript's authors in the intended authorship order]
    Jiaming Wei, Zekun Wu, Adriano Koshiyama, Mar\'ia P\'erez-Ortiz.

    \item[(o) Stage of publication (e.g.\ in submission)]
    Accepted at the REALM Workshop (EMNLP~2026) on 8~September~2026 under that
    venue's non-archival option, camera-ready due 14~September~2026; under
    review at VLM4RWD (NeurIPS~2026). Not published: acceptance at a
    non-archival venue produces no version of record.
\end{description}

\vspace{0.4cm}

\noindent\textbf{3. For multi-authored work, please give a statement of
contribution covering all authors (if single-author, please skip to section 4)}

\begin{description}[leftmargin=1.4em, style=nextline, itemsep=0.3em, topsep=0.3em]
    \item[Jiaming Wei] Designed the study; built the experiment harness and the
    analysis pipeline; ran all reported experiments; performed the analysis;
    produced the figures; wrote the original draft.

    \item[Zekun Wu] Provided the structural writing guidance that shapes how the
    argument separates measured quantities from inferred ones; advised on
    methodology.

    \item[Adriano Koshiyama] Provided the compute infrastructure on which the
    local backbones were run, and the API access through which the
    proxy-served backbone was queried.

    \item[Mar\'ia P\'erez-Ortiz] Supervised the work, and shaped its scope by
    repeatedly redirecting it from what the system could be made to do towards
    what the evidence could actually support.
\end{description}

\noindent All authors reviewed and edited the manuscript.

\vspace{0.5cm}

\noindent\textbf{4. In which chapter(s) of your thesis can this material be
found?}

\vspace{0.2cm}

\noindent
Chapters~\ref{ch:repr}, \ref{ch:pred} and~\ref{ch:e2e} report the same
measurements as the manuscript, and Figures~\ref{fig:overview},
\ref{fig:space} and~\ref{fig:ceilings} are reused from it. The thesis is not an
expansion of the manuscript: Chapters~\ref{ch:intro} and~\ref{ch:bg}, the
benchmark characterisation, and Appendices~\ref{app:ceiling}--\ref{app:repro}
were written for the thesis and have no counterpart in the eight-page workshop
submissions.

\vspace{0.5cm}

\noindent\textbf{5. e-Signatures confirming that the information above is
accurate (this form should be co-signed by the supervisor/senior author unless
this is not appropriate, e.g.\ if the paper was a single-author work)}

\vspace{0.6cm}

\noindent Candidate: \authorName \hfill Date: \submissionDate

\vspace{1.1cm}

\noindent Supervisor/Senior Author: \primarySupervisor \hfill Date:

\vspace{0.4cm}
\thispagestyle{rpdf}
}
